% ============================================================
% DOCUMENTO: Apuntes de Legislación y Práctica Profesional
%             Métodos Alternativos de Resolución de Conflictos
%             y el proceso de desalojo por falta de pago
% AUTOR:      Isaac Edgar Camacho Ocampo
% AÑO:        2026
% ============================================================

% ============================
% CLASE DEL DOCUMENTO
% ============================
\documentclass[12pt,oneside]{book}

% ============================
% CODIFICACIÓN Y LENGUAJE
% ============================
\usepackage[utf8]{inputenc}
\usepackage[T1]{fontenc}
\usepackage[spanish,es-nodecimaldot]{babel}

% ============================
% MÁRGENES
% ============================
\usepackage[
    top=2.5cm,
    bottom=2.5cm,
    left=3cm,
    right=2.5cm,
    headheight=15pt
]{geometry}

% ============================
% MATEMÁTICA
% ============================
\usepackage{amsmath}
\usepackage{amssymb}
\usepackage{mathtools}

% ============================
% IMÁGENES
% ============================
\usepackage{graphicx}
\usepackage{float}
\usepackage{caption}
\usepackage{subcaption}

\graphicspath{
    {./img/}
    {./graficos/}
    {./figuras/}
}

% ============================
% TABLAS
% ============================
\usepackage{booktabs}
\usepackage{array}
\usepackage{longtable}
\usepackage{tabularx}

% ============================
% LISTAS
% ============================
\usepackage{enumitem}

% ============================
% COLORES Y CAJAS
% ============================
\usepackage{xcolor}
\usepackage[most]{tcolorbox}

% ============================
% ENCABEZADOS
% ============================
\usepackage{fancyhdr}

% ============================
% CONTROL TIPOGRÁFICO
% ============================
\usepackage{ragged2e}

% ============================
% HIPERVÍNCULOS Y METADATOS PDF
% ============================
\usepackage[
    colorlinks=true,
    linkcolor=blue!50!black,
    citecolor=green!40!black,
    urlcolor=blue!60!black,
    pdftitle={Apuntes de Legislaci\'on y Pr\'actica Profesional},
    pdfauthor={Isaac Edgar Camacho Ocampo},
    pdfsubject={Apuntes universitarios},
    bookmarks=true
]{hyperref}

% ============================
% METADATOS DEL DOCUMENTO
% ============================
\newcommand{\universidad}{Universidad de Buenos Aires}
\newcommand{\facultad}{Carrera de Arquitectura}
\newcommand{\carrera}{Carrera de Arquitectura}
\newcommand{\materia}{Legislaci\'on y Pr\'actica Profesional}
\newcommand{\titulo}{M\'etodos Alternativos de Resoluci\'on de Conflictos\\[0.3em]y el proceso de desalojo por falta de pago}
\newcommand{\autor}{Isaac Edgar Camacho Ocampo}
\newcommand{\anio}{2026}

% ============================
% CONFIGURACIÓN DE PÁRRAFOS
% ============================
\setlength{\parindent}{0pt}
\setlength{\parskip}{0.5em}

% ============================
% CONFIGURACIÓN DE LISTAS
% ============================
\setlist[itemize]{
    topsep=0.3em,
    itemsep=0.2em
}
\setlist[enumerate]{
    topsep=0.3em,
    itemsep=0.2em
}

% ============================
% CONTROL DE VIUDAS Y HUÉRFANAS
% ============================
\clubpenalty=10000
\widowpenalty=10000

% ============================
% ÍNDICE
% ============================
\setcounter{tocdepth}{2}

% ============================
% COMANDOS MATEMÁTICOS
% ============================
\newcommand{\abs}[1]{\left|#1\right|}
\newcommand{\norm}[1]{\left\lVert#1\right\rVert}
\newcommand{\vect}[1]{\mathbf{#1}}
\newcommand{\deriv}[2]{\frac{d#1}{d#2}}
\newcommand{\pderiv}[2]{\frac{\partial#1}{\partial#2}}

% ============================
% CAJAS ACADÉMICAS
% ============================

% Definición
\newtcolorbox{definition}[1]{
    enhanced,
    breakable,
    title={Definici\'on: #1},
    fonttitle=\bfseries,
    colback=blue!3,
    colframe=blue!50!black,
    boxrule=0.8pt,
    arc=2mm
}

% Importante
\newtcolorbox{important}{
    enhanced,
    breakable,
    title={Importante},
    fonttitle=\bfseries,
    colback=yellow!8,
    colframe=orange!70!black,
    boxrule=0.8pt,
    arc=2mm
}

% Ejemplo
\newtcolorbox{example}[1]{
    enhanced,
    breakable,
    title={Ejemplo: #1},
    fonttitle=\bfseries,
    colback=green!4,
    colframe=green!40!black,
    boxrule=0.8pt,
    arc=2mm
}

% Fórmula / Regla
\newtcolorbox{formula}[1]{
    enhanced,
    breakable,
    title={F\'ormula / Regla: #1},
    fonttitle=\bfseries,
    colback=purple!4,
    colframe=purple!50!black,
    boxrule=0.8pt,
    arc=2mm
}

% Ejercicio
\newtcolorbox{exercise}[1]{
    enhanced,
    breakable,
    title={Ejercicio: #1},
    fonttitle=\bfseries,
    colback=gray!5,
    colframe=gray!60!black,
    boxrule=0.8pt,
    arc=2mm
}

% Nota
\newtcolorbox{note}{
    enhanced,
    breakable,
    title={Nota},
    fonttitle=\bfseries,
    colback=white,
    colframe=gray!50,
    boxrule=0.6pt,
    arc=2mm
}

% Referencia normativa (variante visual)
\newtcolorbox{normref}[1]{
    enhanced,
    breakable,
    title={Referencia normativa: #1},
    fonttitle=\bfseries,
    colback=blue!2,
    colframe=blue!30!black,
    boxrule=0.6pt,
    arc=2mm
}

% ============================
% ENCABEZADOS Y PIES DE PÁGINA
% ============================
\pagestyle{fancy}
\fancyhf{}
\fancyhead[L]{\nouppercase{\leftmark}}
\fancyhead[R]{\materia}
\fancyfoot[C]{\thepage}
\renewcommand{\headrulewidth}{0.4pt}
\renewcommand{\footrulewidth}{0pt}

\fancypagestyle{plain}{
    \fancyhf{}
    \fancyfoot[C]{\thepage}
    \renewcommand{\headrulewidth}{0pt}
}

% ============================================================
% INICIO DEL DOCUMENTO
% ============================================================
\begin{document}

% ============================
% PORTADA
% ============================
\begin{titlepage}

\thispagestyle{empty}

\begin{center}

\vspace*{1cm}

{\large \textsc{\universidad}}

\vspace{0.3cm}

{\large \textsc{\facultad}}

\vspace{3cm}

{\Huge \textbf{\materia}}

\vspace{1cm}

{\LARGE \titulo}

\vspace{0.8cm}

\textit{Comprender el sistema es el primer paso para actuar con criterio frente a cualquier conflicto.}

\vspace{1.2cm}

\rule{0.70\textwidth}{0.5pt}

\vspace{0.5cm}

{\large Apuntes de estudio}

\vspace{0.3cm}

{\normalsize Ciclo lectivo 2026}

\vspace{0.5cm}

\rule{0.70\textwidth}{0.5pt}

\vfill

{\large \autor}

\vspace{0.3cm}

{\large \anio}

\end{center}

\vfill

\begin{flushleft}
\textit{Este es un modesto aporte a los estudiantes de ``Legislaci\'on y Pr\'actica Profesional'' no pretende sustituir el material oficial de la materia.}
\end{flushleft}

\end{titlepage}

% ============================
% ÍNDICE
% ============================
\frontmatter
\tableofcontents
\cleardoublepage

% ============================
% CONTENIDO PRINCIPAL
% ============================
\mainmatter

% ============================================================
\chapter{Fundamentos del sistema de resoluci\'on de conflictos}
% ============================================================

\section{La crisis de la justicia y el rol del derecho procesal}

La percepción generalizada sobre el funcionamiento del sistema judicial
se caracteriza, entre otros aspectos, por su lentitud y por el peso
de la burocracia procesal. Sin embargo, no toda cuestión reviste la
misma complejidad: un conflicto sencillo y uno complejo no deberían
tramitarse de manera idéntica. Por ello, determinados procesos---como
el desalojo por falta de pago---cuentan con un trámite expeditivo y
sustancialmente más ágil que el proceso ordinario.

A partir de esa constatación surge una reflexión central sobre la
relación entre las formas procesales y el conflicto que buscan resolver.
El sistema procesal tiende a operar como un \textbf{molde único} en el
que se encajan todos los conflictos, cuando en rigor debería existir una
diversidad de moldes capaces de dar respuesta adecuada a cada situación
concreta. Sería razonable otorgar al juez mayor flexibilidad para
adecuar las formas al conflicto.

\begin{definition}{Derecho procesal instrumental}
Las formas procesales son una \textbf{garantía} de certeza y seguridad
jurídica, no un fin en sí mismo. El derecho procesal es \textbf{instrumental}
respecto del derecho de fondo: sirve a su realización. Cuando el proceso
judicial, tal como está diseñado, no da una respuesta adecuada al
conflicto, aparecen los \textbf{métodos alternativos de resolución de
conflictos (MASC)} como vía complementaria---no sustitutiva---de la
jurisdicción estatal.
\end{definition}

\begin{important}
La forma es garantía, no finalidad. El derecho de fondo no debe
subordinarse a la forma procesal: si el proceso no permite realizar el
derecho de fondo, no está cumpliendo su función.
\end{important}

\section{Mapa de los sistemas de resoluci\'on de conflictos}

Los distintos mecanismos de resolución de conflictos pueden ordenarse
en dos grandes familias, según la estructura del enfrentamiento entre
las partes y la capacidad del tercero interviniente para imponer una
solución.

\subsection{Sistemas no adversariales}

En los sistemas no adversariales no existe un enfrentamiento directo
entre las partes. Dentro de este grupo se distinguen dos subcategorías
según la presencia o no de un tercero:

\begin{definition}{Autocomposici\'on}
Las propias partes resuelven el conflicto \textbf{sin intervención de
un tercero}. El caso paradigmático es la \textbf{negociación directa}
entre ellas o entre sus representantes legales.
\end{definition}

\begin{definition}{Heterocomposici\'on asistida}
Interviene un tercero que ayuda a las partes a acercar posiciones,
pero \textbf{sin poder imponerles una solución}. Aquí se ubican la
\textbf{mediación} y la \textbf{conciliación}.
\end{definition}

\subsection{Sistemas adversariales}

En los sistemas adversariales existe un enfrentamiento directo entre
las partes. El tercero que interviene \textbf{sí tiene poder para
imponer una solución obligatoria}. Son los casos del \textbf{arbitraje}
y del \textbf{proceso judicial}.

\begin{important}
La diferencia esencial entre los sistemas no adversariales
(negociación, mediación, conciliación) y los adversariales
(arbitraje, proceso judicial) no reside en la presencia de un tercero,
sino en si ese tercero tiene o no \textbf{poder para imponer una
solución obligatoria} a las partes.
\end{important}

El siguiente cuadro sintetiza el mapa de los sistemas:

\begin{table}[H]
\centering
\begin{tabularx}{\textwidth}{
    >{\raggedright\arraybackslash}p{0.22\textwidth}
    >{\raggedright\arraybackslash}p{0.22\textwidth}
    >{\raggedright\arraybackslash}X
    >{\raggedright\arraybackslash}p{0.18\textwidth}
}
\toprule
\textbf{Sistema} & \textbf{Modalidad} & \textbf{Mecanismo} & \textbf{Tercero impone} \\
\midrule
No adversarial & Autocomposición & Negociación directa & No \\
No adversarial & Heterocomposición & Mediación & No \\
No adversarial & Heterocomposición & Conciliación & No \\
Adversarial & Heterocomposición & Arbitraje & Sí \\
Adversarial & Heterocomposición & Proceso judicial & Sí \\
\bottomrule
\end{tabularx}
\caption{Clasificación de los sistemas de resolución de conflictos}
\end{table}

% ============================================================
\chapter{La negociaci\'on como primer paso}
% ============================================================

\section{Los cuatro principios de la negociaci\'on por intereses}

La negociación eficaz entre partes---o entre sus abogados---se rige por
cuatro principios fundamentales del método de negociación basado en
intereses.

\subsection{Primer principio: separar a las personas del problema}

Antes de abordar el conflicto de fondo, deben separarse las cuestiones
personales de las cuestiones sustantivas. Las relaciones personales
deben preservarse: los conflictos interpersonales y los conflictos de
intereses son dimensiones distintas que no deben mezclarse. En el
contexto profesional jurídico, pelearse con un colega cada vez que se
litiga duplica el problema. La recomendación es trabajar sobre los
intereses \textbf{con el apoyo de las personas}, no en contra de ellas.

\subsection{Segundo principio: negociar sobre intereses, no sobre posiciones}

Una \textbf{posición} es la demanda concreta que cada parte sostiene
(``quiero esto'', ``te ofrezco tanto''). Cuando se negocia por
posiciones, las partes ceden de a poco y el resultado suele ser un
punto medio. Detrás de cada posición existe un \textbf{interés real}:
la razón subyacente por la que cada parte quiere lo que pide. Identificar
y trabajar sobre esos intereses---que a veces son comunes o
compatibles---abre posibilidades de acuerdo que la negociación por
posiciones no permite.

\subsection{Tercer principio: generar opciones de beneficio mutuo}

Antes de cerrar una posición definitiva, conviene realizar una
\textbf{tormenta de ideas}: colocar sobre la mesa varias soluciones
posibles y evaluarlas antes de elegir. Adoptar una posición de entrada
es riesgoso porque luego resulta difícil retroceder sin perder
credibilidad.

\subsection{Cuarto principio: usar criterios objetivos}

El cuarto elemento es buscar un \textbf{criterio objetivo} que permita
evaluar las propuestas de manera racional: un valor de mercado, un
antecedente, una tasación, un índice de referencia. Conocer con
claridad el interés real de la otra parte y no apresurarse a cerrar
una posición final son condiciones para una negociación exitosa.

\begin{example}{Las dos empresas y la naranja}
Dos empresas disputaban la misma materia prima, cada una sosteniendo
la posición de quedarse con el 100\% del recurso. Mientras negociaban
desde sus posiciones, no había solución posible.

Al indagar los intereses reales, se descubrió que una empresa necesitaba
la \textbf{cáscara} del producto (para fabricar mermelada) y la otra
necesitaba el \textbf{jugo} (para elaborar jugo). Una vez resuelto el
conflicto sobre la base de los intereses---y no de las posiciones
enfrentadas---ambas partes pudieron quedarse con el 100\% de lo que
efectivamente necesitaban.

\textbf{Conclusión:} cuando se negocia por posiciones, el resultado
es casi siempre una repartición a mitades; cuando se negocia por
intereses, es posible que ambas partes obtengan mucho más.
\end{example}

\section{Vías de hecho prohibidas: el caso del candado}

Para ilustrar los límites legales de la autotutela, resulta útil
analizar un caso práctico frecuente en la consulta profesional: un
locador que, ante la falta de pago del último alquiler, ingresa al
local comercial y le coloca un candado al locatario.

Ante este planteo, la primera reacción no debe ser interponer una
demanda directamente. El primer paso es intentar \textbf{negociar}.
Si la otra parte tiene representación legal, la negociación debe
conducirse con su abogado, no directamente con el cliente. Sin embargo,
más allá de cómo termine la negociación, existe un principio
indiscutible:

\begin{important}
Por más que exista incumplimiento del locatario---incluso falta de
pago---el locador \textbf{no puede recuperar la tenencia del inmueble
por la fuerza}. Las vías de hecho contra quien tiene la tenencia están
expresamente prohibidas. El locador debe necesariamente acudir a la
\textbf{acción de desalojo} por la vía judicial correspondiente.
\end{important}

\begin{normref}{Prohibición de vías de hecho}
El Código Civil y Comercial de la Nación prohíbe al propietario o
locador recuperar la tenencia de un inmueble por propia autoridad,
aun cuando exista incumplimiento del locatario. La tenencia solo puede
recuperarse mediante la vía judicial correspondiente (acción de
desalojo). Actuar por la fuerza expone al locador a una \textbf{denuncia
penal} por usurpación y a la obligación de restituir la tenencia.
\end{normref}

\begin{note}
El cambio de cerradura o la colocación de un candado puede ser
fácilmente objeto de una denuncia penal y resulta sencillo de comprobar.
La situación es diferente en el caso de una \textbf{ocupación ilegal
reciente}: en el momento mismo del hecho, el titular puede actuar en
defensa de la posesión, pero cuanto más tiempo transcurre y más se
consolida la ocupación, menor es el margen para actuar por medios
propios y mayor la exposición a responsabilidad por la propia agresión.
\end{note}

\subsection{Caso especial: desalojo por autoridad administrativa}

Si el inmueble fue desalojado por una autoridad (por ejemplo, clausura
administrativa por riesgo de derrumbe) y ya no hay resistencia por
parte del ocupante, quien era poseedor o dueño puede en principio
recuperarlo, salvo que haya intervenido como parte en el procedimiento
que originó la clausura. En ese caso, debe pedir autorización judicial.
Si no fue parte del procedimiento, puede recuperarlo libremente, siempre
que no exista otro impedimento (proceso penal en curso, problema de
violencia, investigación activa que mantenga afectado el inmueble).

% ============================================================
\chapter{Mediaci\'on y conciliaci\'on}
% ============================================================

\section{Diferencias entre mediaci\'on y conciliaci\'on}

Tanto la mediación como la conciliación son mecanismos de
heterocomposición asistida por un tercero. Sin embargo, presentan una
diferencia funcional esencial.

\begin{definition}{Mediaci\'on}
El \textbf{mediador} asiste a las partes para que sean \textbf{ellas
mismas} quienes encuentren la solución al conflicto. No puede proponer
fórmulas conciliatorias ni resolver el caso por sí mismo.
\end{definition}

\begin{definition}{Conciliaci\'on}
El \textbf{conciliador} puede proponer \textbf{fórmulas conciliatorias}
a las partes. Sin embargo, tampoco puede resolver el caso por sí mismo:
no decide el conflicto.
\end{definition}

\begin{table}[H]
\centering
\begin{tabularx}{\textwidth}{
    >{\raggedright\arraybackslash}X
    >{\centering\arraybackslash}p{0.22\textwidth}
    >{\centering\arraybackslash}p{0.22\textwidth}
}
\toprule
\textbf{Característica} & \textbf{Mediador} & \textbf{Conciliador} \\
\midrule
Asiste a las partes & Sí & Sí \\
Puede proponer fórmulas & No & Sí \\
Puede imponer una solución & No & No \\
Decide el conflicto & No & No \\
\bottomrule
\end{tabularx}
\caption{Comparación entre mediador y conciliador}
\end{table}

\subsection{El juez como conciliador: la audiencia del art.\ 360 CPCCN}

En la \textbf{audiencia preliminar} regulada por el artículo 360 del
Código Procesal Civil y Comercial de la Nación (CPCCN), el juez actúa
con facultades de conciliador: puede proponer fórmulas conciliatorias
a las partes. Si la conciliación no prospera, el juez fija los
\textbf{hechos controvertidos} del proceso y ordena los medios de prueba
que se producirán.

\begin{normref}{Audiencia preliminar (art.\ 360 CPCCN)}
El artículo 360 del CPCCN regula la audiencia preliminar: el juez
invita a las partes a una conciliación y, si esta no prospera, fija
los hechos controvertidos, resuelve las cuestiones relativas a la
prueba ofrecida y ordena las medidas necesarias para el desarrollo
del proceso.
\end{normref}

\begin{note}
El propio código resuelve la tensión entre proponer fórmulas
conciliatorias y el riesgo de \textbf{prejuzgamiento}, estableciendo
expresamente que el juez no incurre en prejuzgamiento por formular
propuestas conciliatorias en esa audiencia. Se trata de una facultad
habilitada para favorecer el acuerdo, distinta de anticipar el sentido
del fallo.
\end{note}

En esta audiencia, conviene que las partes hablen directamente---no
solo sus abogados---dado que la mediación y la conciliación presuponen
que sean las propias partes quienes busquen la solución. Ello pone de
relieve la importancia de la \textbf{confidencialidad}.

\begin{definition}{Principio de confidencialidad}
Todo lo manifestado por las partes durante una audiencia de mediación
o conciliación es \textbf{confidencial} y no puede ser utilizado como
prueba en el proceso judicial posterior si no se llega a un acuerdo.
El mediador tampoco puede ser citado como testigo respecto de lo
tratado en la audiencia. Esta garantía es la que permite a las partes
hablar con franqueza sobre los hechos sin temor a que ello las
perjudique en una etapa ulterior.
\end{definition}

\section{Procedimiento, plazos y confidencialidad de la mediaci\'on}

\subsection{Tipos de mediación}

La mediación prejudicial puede ser \textbf{privada} u \textbf{oficial}.

\begin{itemize}
    \item \textbf{Mediación privada:} la parte requirente elige un mediador
    de la lista habilitada por el Ministerio de Justicia y debe proponerle
    \textbf{tres mediadores} a la otra parte para que el requerido elija
    uno de ellos.
    \item \textbf{Mediación oficial:} se presenta un formulario ante la
    mesa de entradas de la cámara que corresponda---actualmente por vía
    electrónica---y el mediador oficial es designado \textbf{por sorteo}.
\end{itemize}

\subsection{Plazos y efectos del acta de cierre}

La mediación tiene un plazo de \textbf{60 días} para el desarrollo de
las audiencias, prorrogable. Rige la confidencialidad. Si no se alcanza
un acuerdo, se labra un \textbf{acta de cierre de la mediación}, que
debe acompañarse como \textbf{requisito de admisibilidad} de la demanda:
si el proceso judicial se inicia sin haberse transitado la mediación
previa, el juez remite el expediente a mediación antes de continuar.

\subsection{Razones de la mediación previa obligatoria}

La mediación prejudicial obligatoria fue instituida porque el proceso
judicial se encontraba congestionado y no brindaba una respuesta
oportuna. Dentro del proceso judicial ya iniciado, existe además una
segunda oportunidad de conciliación: la audiencia del artículo 360
CPCCN.

\begin{note}
Una diferencia clave entre la mediación y la audiencia preliminar es
el \textbf{principio de informalidad}: en la mediación no existe
obligación de exhibir la prueba (las partes pueden mencionar que
cuentan con determinados elementos sin mostrarlos), mientras que en
la audiencia preliminar ya se ha formulado la demanda, se ha fundado
la pretensión y se ha ofrecido la prueba---toda la evidencia está
``sobre la mesa''. Además, en esa segunda instancia el juez ya conoce
la demanda y puede proponer fórmulas conciliatorias con mayor
conocimiento del caso.
\end{note}

\subsection{Sanci\'on por incomparecencia}

\begin{normref}{Evolución normativa: sanción por incomparecencia a la mediación}
El régimen original de la Ley de Mediación Prejudicial Obligatoria
preveía la imposición \textbf{obligatoria} de una multa a la parte
requerida que no comparecía sin causa justificada. La reforma
introducida por la \textbf{Ley 26.589} sustituyó ese mandato imperativo
por una \textbf{facultad del juez} (``podrá'' imponer la multa), lo
que en la práctica diluyó considerablemente su aplicación.
\end{normref}

La consecuencia práctica de ese cambio de redacción es que casi no se
aplican multas. Una multa disuasiva sería muy efectiva frente a quienes
firman el acuerdo de comparecencia y luego no concurren a las audiencias.

\section{Mediaci\'on, prescripci\'on y caducidad}

La \textbf{prescripción} es el instituto por el cual se extingue la
acción ante la inacción del titular durante el plazo establecido por
la ley. La prescripción puede interrumpirse o suspenderse.

\begin{definition}{Efecto de la mediaci\'on sobre la prescripci\'on}
Cuando se inicia el trámite de mediación, se \textbf{suspende el curso
de la prescripción}. La suspensión se extiende hasta \textbf{veinte
días después} de labrada el acta de cierre de la mediación, momento
a partir del cual el plazo remanente vuelve a correr.
\end{definition}

\begin{example}{C\'alculo del plazo remanente}
Supóngase que a una acción le quedaban solo 7 días de plazo de
prescripción cuando se inició la mediación. La mediación dura,
supóngase, 90 días. Cuando termina la mediación y se labra el acta
de cierre, transcurren 20 días más. Recién entonces comienza a correr
nuevamente el plazo: en este ejemplo, solo \textbf{7 días}.

Si no se inicia la demanda dentro de esos 7 días, la acción
\textbf{prescribe} y ya no puede interponerse la demanda sin exponerse
a la excepción de prescripción.
\end{example}

\begin{important}
Cuando el cliente manifiesta que ``el deudor es de su absoluta
confianza y jamás dejará prescribir la acción'', el abogado debe
advertirle expresamente---y \textbf{por escrito}, mediante un
consentimiento informado---que el plazo está corriendo y que dejarlo
vencer implica perder la acción, con independencia de cualquier
consideración personal. La práctica ilustra que esta advertencia
es esencial: la confianza del cliente en su deudor no es oponible
frente a la excepción de prescripción.
\end{important}

La misma lógica de plazos se aplica a las \textbf{medidas cautelares},
como se desarrollará en el capítulo siguiente.

% ============================================================
\chapter{Arbitraje}
% ============================================================

\section{Cl\'ausula compromisoria y compromiso arbitral}

El arbitraje es el segundo mecanismo adversarial. A diferencia del
proceso judicial, en el arbitraje son las propias partes quienes se
\textbf{sustraen voluntariamente de la jurisdicción estatal} y someten
su conflicto a la decisión de uno o más árbitros.

\begin{definition}{Cl\'ausula compromisoria}
Pacto incluido en el contrato \textbf{antes de que exista el conflicto},
por el cual las partes acuerdan que las controversias futuras derivadas
de ese contrato serán resueltas por arbitraje. Puede determinar el
árbitro concreto o remitir al procedimiento estandarizado de una
institución (por ejemplo, una cámara de comercio).
\end{definition}

\begin{definition}{Compromiso arbitral}
Instrumento suscripto \textbf{una vez que el conflicto ya existe}, que
fija en concreto las cuestiones específicas que el árbitro deberá
resolver. Los árbitros no pueden fallar sobre cuestiones no sometidas
al compromiso ni fuera del plazo allí establecido.
\end{definition}

\subsection{Funciones de la cl\'ausula compromisoria}

La cláusula compromisoria cumple dos funciones:

\begin{itemize}
    \item \textbf{Función positiva:} si surge el conflicto y la otra
    parte no quiere firmar el compromiso arbitral, puede demandarse
    judicialmente que se constituya ese compromiso; si la parte no
    comparece, el propio juez puede suplir esa voluntad y designar
    el árbitro.
    \item \textbf{Función negativa:} si pese a la existencia de la
    cláusula compromisoria se interpone una demanda ante la justicia
    ordinaria, la parte demandada puede oponer la excepción de
    incompetencia, ya que el competente es el tribunal arbitral.
\end{itemize}

\subsection{Jurisdicci\'on, competencia y pr\'orroga}

\begin{definition}{Jurisdicci\'on y competencia}
La \textbf{jurisdicción} es la potestad de juzgar, atribuida a todos
los jueces. La \textbf{competencia} es la \textit{medida} de esa
jurisdicción para un caso concreto: se divide por materia, territorio
y personas.
\end{definition}

Solo puede prorrogarse la competencia \textbf{territorial} y en asuntos
\textbf{patrimoniales}. No es posible prorrogar la competencia por
razón de la materia ni por razón de las personas.

\begin{example}{Pr\'orroga de competencia en un contrato de locaci\'on}
Un contrato de locación es una cuestión patrimonial (cesión de la
tenencia de un inmueble a cambio de un precio en dinero). Por eso es
válido prorrogar la competencia territorial.

Supóngase un inmueble ubicado en Catamarca, cuyo contrato se firmó en
Córdoba, con una cláusula que fija la competencia de los tribunales
de la Ciudad de Buenos Aires. Si la demanda de desalojo se interpone
en Córdoba, la otra parte puede oponer \textbf{excepción de
incompetencia} por incumplimiento de la cláusula de prórroga, y el
planteo prosperará.
\end{example}

\begin{note}
La competencia territorial y patrimonial es \textbf{disponible} para
las partes: por eso puede prorrogarse y, en principio, el juez no
puede declararla de oficio. Distinto es el caso de la competencia
federal frente a la local: si se demanda en el fuero equivocado, el
juez puede declararse incompetente de oficio, ya que no es una materia
disponible para las partes.

Del mismo modo, si existe una cláusula compromisoria y se inicia
demanda por vía judicial, la otra parte puede oponer la excepción
correspondiente, que prosperará.
\end{note}

\section{Los l\'imites del \'arbitro frente al juez}

\subsection{Funci\'on jurisdiccional del \'arbitro}

Se discute doctrinariamente si la función del árbitro constituye o no
una función jurisdiccional en sentido estricto. Si se entiende la
jurisdicción como la función de \textbf{fijar los hechos
controvertidos, determinar si existe una norma que los contemple como
antecedente y elevar esos hechos a la norma para dictar una decisión
concreta}, entonces el árbitro cumple exactamente esa función: fija los
hechos, los subsume en una norma y dicta un \textbf{laudo}, que cumple
la misma función que la sentencia judicial.

\subsection{El \'arbitro carece de imperium}

\begin{definition}{Imperium}
Es el \textbf{poder coactivo} del Estado que permite al juez ejecutar
sus decisiones de manera forzada, incluyendo el empleo de la fuerza
pública.
\end{definition}

\begin{important}
El árbitro \textbf{carece de imperium}: no puede ejecutar coactivamente
su propio laudo. Si la parte condenada no cumple voluntariamente la
decisión arbitral, la parte interesada debe acudir al \textbf{juez
estatal que hubiera sido competente} de no existir la cláusula arbitral,
para que sea él quien ordene la ejecución forzada. Tampoco puede el
árbitro trabar medidas cautelares por sí mismo.
\end{important}

Esta limitación tiene una consecuencia práctica muy relevante: en
conflictos donde lo prioritario es recuperar un inmueble cuanto antes
(como en el desalojo), el arbitraje puede suponer un paso adicional---
obtener el laudo y luego acudir al juez para que ordene el
lanzamiento---lo que en algunos casos puede demandar más tiempo que
recurrir directamente a la vía judicial.

\subsection{Diferencias de fundamentaci\'on: derecho vs.\ equidad}

Mientras el derecho positivo exige que las decisiones judiciales estén
fundadas conforme a derecho, un \textbf{arbitraje de amigables
componedores} puede resolver según la \textbf{equidad}, sin la misma
exigencia de fundamentación normativa exhaustiva, resolviendo ``a
verdad sabida y buena fe guardada''.

\subsection{Pericia arbitral}

Cuando el conflicto requiere un conocimiento técnico específico---por
ejemplo, liquidar una sociedad compleja---puede designarse un
\textbf{perito} (o un cuerpo de peritos) para que decida esa cuestión
técnica, con carácter \textbf{vinculante} para las partes. Esta figura
se denomina \textbf{pericia arbitral}.

% ============================================================
\chapter{Aplicaci\'on pr\'actica: el caso de desalojo}
% ============================================================

\section{Medidas cautelares y su caducidad}

\begin{definition}{Medidas cautelares}
Las medidas cautelares son \textbf{provisionales} y \textbf{accesorias}
a un proceso principal. No constituyen un proceso en sí mismas: su
función es proteger el resultado eventual del proceso principal
mientras este se sustancia (por ejemplo, un embargo o un secuestro de
mercadería).
\end{definition}

\begin{normref}{Caducidad de las medidas cautelares}
Por su carácter accesorio, la ley fija un \textbf{plazo de caducidad}:
cuando una medida cautelar se traba antes de la interposición de la
demanda, quien la obtuvo debe iniciar el proceso principal dentro del
plazo de \textbf{diez días}, contados desde que la medida quedó
trabada. Vencido ese plazo sin haberse iniciado la demanda, la medida
cautelar caduca de pleno derecho, y quien la solicitó puede quedar
obligado a responder por los daños y perjuicios causados por la traba
injustificada.
\end{normref}

\subsection{Interacci\'on entre cautelar y mediaci\'on obligatoria}

Si para iniciar la demanda resulta necesario transitar antes la
mediación obligatoria, el esquema opera del siguiente modo: si la
mediación se inicia \textbf{dentro} del plazo de 10 días, ese plazo
se detiene mientras dura el trámite de mediación y vuelve a correr
recién a los \textbf{20 días del acta de cierre}, exactamente igual
que sucede con la prescripción.

\begin{important}
La estrategia óptima es iniciar la mediación en los \textbf{primeros
días} del plazo de 10 días, no en los últimos. Si se inicia el día 9,
quedará apenas 1 día hábil tras el acta de cierre para redactar e
interponer la demanda. Si se inicia el primer día, quedarán
prácticamente los 10 días completos disponibles luego de finalizada la
mediación.
\end{important}

\begin{example}{El boleto de compraventa y la multa diaria}
Un boleto de compraventa incluía una cláusula penal: una multa por
cada día de demora en escriturar. Si el titular de la cautelar la
deja caducar por no iniciar a tiempo la demanda, deberá responder
por esa demora, ya que dejar caducar la medida hace presumir que no
existía razón suficiente para haberla trabado.

Si, en cambio, la demanda se inició en tiempo y forma, no hay
responsabilidad por la demora, porque el reclamo estaba fundado y
la caducidad no llegó a producirse.
\end{example}

\begin{note}
Debe notificarse al mediador el resultado del proceso, ya que sus
honorarios se calculan sobre el monto reclamado y son pagados---según
el caso---por la parte requerida si hay acuerdo, o por quien resulte
condenado en costas si el conflicto termina en juicio.
\end{note}

\section{Resoluci\'on del caso: falta de pago y mejoras no permitidas}

\subsection{Planteo del caso}

\begin{table}[H]
\centering
\begin{tabularx}{\textwidth}{
    >{\raggedright\arraybackslash}p{0.30\textwidth}
    >{\raggedright\arraybackslash}X
}
\toprule
\textbf{Elemento} & \textbf{Descripción} \\
\midrule
Tipo de contrato & Locación (contrato de locación de inmueble) \\
Plazo & 10 años \\
Precio & Pactado con actualización mensual conforme a una tasa de referencia \\
Condición contractual & El contrato prohíbe expresamente al locatario realizar mejoras en el inmueble \\
Incumplimiento 1 & El locatario realizó mejoras no permitidas por el contrato \\
Incumplimiento 2 & El locatario adeuda el pago del alquiler conforme a la forma pactada \\
Consulta & El locador consulta cómo proceder frente a ambos incumplimientos \\
\bottomrule
\end{tabularx}
\caption{Datos del caso práctico de desalojo}
\end{table}

\subsection{Primer an\'alisis: la ley aplicable}

Lo primero que debe analizarse al resolver el caso es
\textbf{qué ley estaba vigente cuando se firmó el contrato}.

\begin{normref}{Irretroactividad de la ley en materia contractual (art.\ 7 CCCN)}
El artículo 7 del Código Civil y Comercial de la Nación establece que
las leyes se aplican a partir de su entrada en vigencia y \textbf{no
tienen efecto retroactivo}, salvo disposición en contrario. En materia
de contratos en curso de ejecución, las nuevas leyes supletorias no se
aplican a las relaciones y situaciones jurídicas ya constituidas, salvo
que la nueva ley sea más favorable al consumidor.
\end{normref}

La ley vigente al momento en que se acordó el contrato es la que
delimita las posibilidades de acción según el incumplimiento. Si el
locador pretendiera aplicar retroactivamente una ley posterior, no
tendría razón.

\subsubsection{Mapa normativo de las locaciones urbanas}

Las locaciones urbanas en Argentina transitaron distintos regímenes
normativos relevantes:

\begin{itemize}
    \item \textbf{Ley 27.551:} introdujo, entre otras disposiciones, la
    actualización \textbf{semestral} por un índice legal, en lugar de
    la actualización mensual libremente pactada.
    \item \textbf{DNU 70/2023:} liberalizó nuevamente la posibilidad de
    pactar libremente el método de actualización, incluso la actualización
    mensual.
\end{itemize}

\begin{important}
La actualización mensual del caso solo sería válida si el contrato se
firmó \textbf{después del DNU 70/2023}. Si se hubiera firmado durante
la vigencia de la Ley 27.551, la actualización mensual libremente
pactada habría sido \textbf{nula} por tratarse de una norma de orden
público, y correspondería aplicar la actualización semestral por el
índice legal previsto en esa ley.
\end{important}

\subsection{La falta de pago: intimaci\'on previa}

\begin{normref}{Requisitos de la intimación de pago previa al desalojo}
Ante la falta de pago del alquiler, el locador debe \textbf{intimar
fehacientemente} al locatario a que abone la suma debida en un plazo
\textbf{no menor a diez días}, especificando el \textbf{monto exacto}
reclamado. Solo vencido ese plazo sin que se regularice el pago, y
encontrándose adeudados al menos \textbf{dos períodos locativos}, queda
expedita la acción de desalojo por falta de pago.
\end{normref}

\subsection{Las mejoras no permitidas}

El Código Civil y Comercial de la Nación establece como causal de
desalojo, además de la falta de pago, el incumplimiento de la
obligación de \textbf{conservar la cosa en el estado en que fue
entregada}. La regla general es que, si el contrato no prohíbe las
mejoras, el locatario puede realizarlas. En el caso analizado, el
contrato las prohibía expresamente, de modo que la mejora realizada
constituye un \textbf{incumplimiento contractual autónomo}.

\subsubsection{Prueba de las mejoras no autorizadas}

Los medios de prueba relevantes para acreditar la existencia de mejoras
no autorizadas son:

\begin{itemize}
    \item \textbf{Inventario del estado inicial del inmueble:} documento
    que describe el estado del inmueble al momento de la entrega al
    locatario.
    \item \textbf{Fotografías:} registro gráfico del estado inicial y del
    estado posterior a la mejora.
    \item \textbf{Acta de constatación notarial:} si es necesario constatar
    formalmente la mejora una vez detectada.
\end{itemize}

\begin{important}
Para que el acta de constatación tenga \textbf{valor probatorio pleno},
debe citarse a la \textbf{parte contraria} a presenciarla. Si no
concurre, conviene llevar testigos, porque de lo contrario el acta
es una prueba unilateral---realizada sin contradicción---y pierde
fuerza frente al juez.
\end{important}

\subsubsection{Concepto de ``mejora'' y sus l\'imites}

Existe una distinción relevante entre una mejora en sentido propio y
los actos de conservación ordinaria del inmueble:

\begin{itemize}
    \item \textbf{Acto de conservación ordinaria:} en un contrato de
    10 años, pintar la casa durante la locación constituye conservación
    normal del inmueble, no una mejora en sentido propio. El contrato
    puede exigir que la devolución se realice en condiciones similares
    a las de la entrega (recién pintada), pero eso no equivale a prohibir
    pintar durante la locación.
    \item \textbf{Mejora en sentido propio:} tirar una pared para ampliar
    un ambiente, instalar una cocina nueva u obras de alteración
    estructural. Estas sí quedan alcanzadas por una prohibición
    contractual de mejoras.
\end{itemize}

\begin{note}
El escribano puede constatar lo que percibe directamente, pero no
siempre es la persona idónea para calificar técnicamente un daño
estructural o edilicio. Según la complejidad del caso, puede requerirse
la intervención de un \textbf{arquitecto o ingeniero} que complemente
la constatación notarial con una opinión técnica fundada.
\end{note}

\subsection{Estrategia: mediaci\'on, arbitraje o proceso judicial}

Recuperar el inmueble cuanto antes resulta siempre conveniente desde
el punto de vista económico: tener a alguien habitando el inmueble sin
pagar y sin cuidarlo implica un deterioro considerable y un significativo
costo de oportunidad para el propietario. Por eso, en este tipo de
conflictos, la elección de la vía de resolución debe hacerse con
criterio estratégico.

\begin{table}[H]
\centering
\begin{tabularx}{\textwidth}{
    >{\raggedright\arraybackslash}p{0.20\textwidth}
    >{\raggedright\arraybackslash}X
    >{\raggedright\arraybackslash}p{0.18\textwidth}
}
\toprule
\textbf{Vía} & \textbf{Características principales} & \textbf{Costo / tiempo} \\
\midrule
Proceso judicial & Garantía de una decisión ejecutable directamente por el propio juez, incluyendo el lanzamiento. Mayor costo y mayor tiempo. & Mayor \\
\addlinespace
Arbitraje & Puede acortar el tiempo de resolución del fondo, pero exige un paso adicional---acudir a la justicia ordinaria---para ejecutar el lanzamiento, ya que el árbitro carece de imperium. & Intermedio \\
\addlinespace
Mediación & Permite un acuerdo voluntario ejecutable sin necesidad de homologación (salvo intereses de menores o incapaces) y evita años de proceso judicial. Opción generalmente más eficiente en tiempo y costo para este tipo de conflictos. & Menor \\
\bottomrule
\end{tabularx}
\caption{Comparación de vías de resolución en conflictos de desalojo}
\end{table}

\subsubsection{Homologaci\'on del acuerdo de mediaci\'on}

\begin{note}
En el arbitraje, si el laudo no se cumple voluntariamente, debe
acudirse al juez competente para que ordene su ejecución.

En la mediación, \textbf{no existe homologación}, salvo que estén
involucrados menores o incapaces (en ese caso sí se exige homologación
judicial). Fuera de esa excepción, el acuerdo de mediación se ejecuta
directamente si no se cumple, con la sola acta firmada por el mediador
como título ejecutivo.
\end{note}

\subsection{Conclusi\'on del caso}

El análisis del caso involucra:

\begin{enumerate}
    \item Identificar la \textbf{ley aplicable} según la fecha del contrato.
    \item Verificar los \textbf{requisitos formales de la intimación de
    pago} (monto exacto, plazo de 10 días, dos períodos adeudados).
    \item Reunir la \textbf{prueba de las mejoras no permitidas}
    (inventario inicial, fotografías, acta de constatación con citación
    de la parte contraria o testigos).
    \item Decidir con criterio profesional \textbf{qué estrategia}---
    negociación, mediación o vía judicial---conviene según el objetivo
    real del cliente: no solo ganar el juicio, sino \textbf{recuperar
    el inmueble en el menor tiempo y al menor costo posible}.
\end{enumerate}

% ============================================================
\chapter{Cuadro Cornell de repaso}
% ============================================================

El siguiente cuadro Cornell sintetiza los conceptos fundamentales
desarrollados a lo largo de los bloques temáticos. La columna izquierda
contiene las preguntas o conceptos clave; la columna derecha, el
desarrollo correspondiente; y la fila final, una síntesis integradora.

\begin{table}[H]
\centering
\begin{tabularx}{\textwidth}{
    >{\raggedright\arraybackslash}p{0.28\textwidth}
    >{\raggedright\arraybackslash}X
}
\toprule
\textbf{Preguntas / palabras clave} & \textbf{Notas / respuestas} \\
\midrule

\textbf{¿Por qué el derecho procesal es ``instrumental''?} &
Porque su función es servir al derecho de fondo, no imponerse
sobre él. Las formas procesales garantizan certeza y seguridad
jurídica, pero cuando el molde procesal único no da una respuesta
adecuada a un conflicto concreto, aparecen los métodos alternativos
de resolución de conflictos (MASC) como complemento del proceso
judicial. \\[0.4em]

\textbf{Sistemas adversariales vs.\ no adversariales} &
En los sistemas no adversariales (negociación, mediación,
conciliación) el tercero---si lo hay---no puede imponer una solución.
En los sistemas adversariales (arbitraje, proceso judicial) sí existe
un tercero con poder para imponer una decisión obligatoria a las
partes. \\[0.4em]

\textbf{Los cuatro principios de la negociación por intereses} &
1) Separar a las personas del problema. 2) Negociar sobre intereses,
no sobre posiciones. 3) Generar opciones de beneficio mutuo antes de
decidir. 4) Usar criterios objetivos (valor de mercado, antecedentes,
tasaciones) para evaluar las propuestas. \\[0.4em]

\textbf{¿Puede el locador recuperar el inmueble por la fuerza si no le pagan?} &
No. Aunque exista falta de pago, está prohibido usar vías de hecho
para recuperar la tenencia. Debe iniciarse la acción de desalojo
correspondiente; actuar por la fuerza expone al locador a una denuncia
penal y a la obligación de restituir la tenencia. \\[0.4em]

\textbf{Mediador vs.\ conciliador} &
El mediador asiste a las partes para que sean ellas quienes encuentren
la solución, sin poder proponer fórmulas. El conciliador sí puede
proponer fórmulas conciliatorias (por ejemplo, el juez en la audiencia
del art.\ 360 CPCCN), pero tampoco puede resolver el caso por sí
mismo. \\[0.4em]

\textbf{¿Cómo afecta la mediación a la prescripción?} &
El inicio de la mediación \textbf{suspende} el curso de la
prescripción. Esa suspensión se extiende hasta 20 días después del
acta de cierre de la mediación, momento en que vuelve a correr el
plazo remanente. Si no se demanda a tiempo dentro de ese plazo, la
acción prescribe. \\[0.4em]

\textbf{Cláusula compromisoria vs.\ compromiso arbitral} &
La cláusula compromisoria se pacta antes del conflicto (habitualmente
en el contrato) y solo prevé que las controversias futuras se
resolverán por arbitraje. El compromiso arbitral se firma una vez
que el conflicto ya existe y fija en concreto la materia que el
árbitro deberá resolver. \\[0.4em]

\textbf{¿Por qué el árbitro necesita al juez para ejecutar su laudo?} &
Porque el árbitro carece de \textbf{imperium}: no tiene poder coactivo
propio. Aunque cumple una función equivalente a la jurisdiccional (fija
hechos, aplica derecho y dicta un laudo), si la parte condenada no
cumple voluntariamente, hay que recurrir al juez estatal competente
para ejecutar la decisión. Tampoco puede trabar medidas cautelares
por sí mismo. \\[0.4em]

\textbf{Caducidad de las medidas cautelares} &
Las medidas cautelares son provisionales y accesorias a un proceso
principal. Si se traban antes de la demanda, hay un plazo de
\textbf{10 días} para iniciar el proceso principal; si no se cumple,
la cautelar caduca y quien la obtuvo puede responder por los daños
causados por la traba injustificada. \\[0.4em]

\textbf{¿Qué ley se aplica a un contrato firmado antes de un cambio legislativo?} &
La ley vigente al momento en que se acordó el contrato, por el
principio de irretroactividad de la ley (art.\ 7 CCCN). Los cambios
legislativos posteriores (por ejemplo, entre la Ley 27.551 y el DNU
70/2023) no modifican retroactivamente las condiciones ya pactadas. \\[0.4em]

\textbf{Requisitos para demandar el desalojo por falta de pago} &
Intimar fehacientemente al locatario a pagar en un plazo no menor
a 10 días, indicando el monto exacto adeudado, y que se adeuden al
menos \textbf{dos períodos} de alquiler. Solo vencido ese plazo sin
regularización queda expedita la acción de desalojo. \\[0.4em]

\midrule

\multicolumn{2}{p{\textwidth}}{
\textbf{Síntesis final:}
El material recorrió el conjunto de herramientas con que cuenta un
profesional para resolver un conflicto contractual, ordenándolas según
el grado de intervención de un tercero y el poder de ese tercero para
imponer una solución: desde la negociación directa---regida por los
cuatro principios de separar personas de problemas, negociar por
intereses, generar opciones de beneficio mutuo y usar criterios
objetivos---pasando por la mediación y la conciliación, donde un
tercero asiste pero no decide, hasta el arbitraje y el proceso
judicial, donde un tercero sí impone una solución obligatoria, con
la diferencia de que solo el juez tiene el imperium necesario para
ejecutarla por la fuerza pública. Todo ese recorrido conceptual se
aplicó a un caso concreto: el desalojo de un inmueble por falta de
pago y por mejoras no permitidas, donde se combinaron reglas de
derecho de fondo (irretroactividad de la ley, requisitos de la
intimación de pago, prohibición de vías de hecho, prueba de las
mejoras) con una decisión estratégica profesional: cuál de todos los
caminos disponibles---negociación, mediación, arbitraje o proceso
judicial---permite recuperar el inmueble en el menor tiempo y al
menor costo posible.
} \\

\bottomrule
\end{tabularx}
\caption*{Cuadro Cornell de repaso integral}
\end{table}

% ============================================================
% FIN DEL DOCUMENTO
% ============================================================
\end{document}
